\chapter{通常のスペクトル対象。}
\VerdierPhysicalPageAnchor{0224}%
\VerdierSemanticAnchor{ch4d-4}
\begingroup
\small
\setlength{\parskip}{0pt}
\setlength{\emergencystretch}{1.5em}
\setlength{\abovedisplayskip}{4pt plus 2pt minus 2pt}%
\setlength{\belowdisplayskip}{4pt plus 2pt minus 2pt}%

\VerdierSemanticAnchor{ch4d-4-1}
\section{第一標準スペクトル対象。}

\VerdierSemanticAnchor{ch4d-4-1-1}
\noindent\textbf{4.1.1.}\quad
$\widetilde{\mathbb Z}$ は順序集合
$\mathbb Z\cup\{-\infty,+\infty\}$、またはこの順序集合に付随する圏
（第 II 章、4.3.1）を表すことを思い出そう。$\mathcal A$ を加法圏とし、
$Y\colon\widetilde{\mathbb Z}\longrightarrow
\operatorname{comp}(\mathcal A)$ を関手とする：

\VerdierSemanticAnchor{ch4d-4-1-1-1}
\begin{equation}
\begin{aligned}
q&\in\widetilde{\mathbb Z},
  &q&\longmapsto Y(q),\\
(p,q)&\in\mathcal{F}\ell_1(\widetilde{\mathbb Z}),
  &u(p,q)&\colon Y(p)\longrightarrow Y(q).
\end{aligned}
\tag{4.1.1.1}
\end{equation}

第 I 章 (3.3.6) では、関手 $Y$ に、$K(\mathcal A)$ に値をとる
（$\widetilde{\mathbb Z}$ 型の）スペクトル対象（第 II 章、4.1.2 と
4.1.3）を対応させた。その構成をここで思い出す。

まず、$\mathcal{F}\ell_1(\widetilde{\mathbb Z})$ の各元 $(p,q)$
（すなわち $p\leq q$ を満たす各組
$p,q\in\widetilde{\mathbb Z}$）に対し、$\mathcal Y(p,q)$ を射
$u(p,q)$ の写像錐（第 I 章、3.1.2）とする：

\VerdierSemanticAnchor{ch4d-4-1-1-2}
\begin{equation}
\begin{aligned}
\mathcal Y(p,q)^i&=Y(q)^i\oplus Y(p)^{i+1},\\
d_{\mathcal Y(p,q)}^i
&=
\begin{pmatrix}
d_{Y(q)}^i & u(p,q)^{i+1}\\
0 & -d_{Y(p)}^{i+1}
\end{pmatrix}.
\end{aligned}
\tag{4.1.1.2}
\end{equation}

$\mathcal{F}\ell_1(\widetilde{\mathbb Z})$ の元 $(p,q)$ と $(p',q')$ の
任意の組で $p\leq p'$ かつ $q\leq q'$ を満たすものに対し、

\VerdierSemanticAnchor{ch4d-4-1-1-3}
\begin{equation}
\Phi\bigl((p,q),(p',q')\bigr)\colon
\mathcal Y(p,q)\longrightarrow\mathcal Y(p',q'),
\tag{4.1.1.3}
\end{equation}
を、次の複体の射のホモトピー類として定義する：

\VerdierSemanticAnchor{ch4d-4-1-1-4}
\begin{equation}
\Psi\bigl((p,q),(p',q')\bigr)\colon
\mathcal Y(p,q)\longrightarrow\mathcal Y(p',q'),
\tag{4.1.1.4}
\end{equation}
この射は次の式で定義される：

\VerdierSemanticAnchor{ch4d-4-1-1-5}
\begin{equation}
\Psi\bigl((p,q),(p',q')\bigr)^i
=
\begin{pmatrix}
u(q,q')^i & 0\\
0 & u(p,p')^{i+1}
\end{pmatrix}.
\tag{4.1.1.5}
\end{equation}
\endgroup
